\textit{Pour chacune des affirmations suivantes, indiquer si elle est \textbf{vraie} ou \textbf{fausse}, en \textbf{justifiant} la réponse. Un réponse non justifiée n'est pas prise en compte.}

\begin{center}
	\textit{Les trois questions sont indépendantes.}
\end{center}

\begin{enumerate}[itemsep=5mm]
	\item On considère un réel $u$.

	On considère sur $\mathbb{R}$ l'équation
	\[(E)\quad x^2+x-u^2=0\]

	\textbf{Affirmation} : Quelle que soit la valeur du réel $u$, l'équation (\textit{E}) possède deux solutions réelles distinctes.

	\item On considère la suite \mbox{$(u_n)$} définie pour tout entier naturel $n$ par
	\[u_n=2^{-n}\]

	\textbf{Affirmation} : La suite \mbox{$(u_n)$} est une suite géométrique de raison \mbox{$\dfrac12$}.

	\item On considère la fonction $f$ définie sur $\mathbb{R}$ par
	\[f(x)=\e^x-1\]

	On note \mbox{$\mathcal{C}$} la courbe représentative de la fonction $f$ dans un repère.

	On note \mbox{$\mathcal{T}$} la tangente à la courbe \mbox{$\mathcal{C}$} au point d'abscisse $0$.

	On note A le point de coordonnées \mbox{$(3;3)$.}

	\textbf{Affirmation} : Le point A appartient à la tangente \mbox{$\mathcal{T}$}.
\end{enumerate}
